\section{制約の順守の判定}
\label{sec_compliance}
\subsection{カルタ}
カルタの制約は，ボケの書き出しのかなと札のかなの一致で判定できる．
札のかなはお題のラベルに含まれないため，同じお題のボケの書き出しのかなの最頻値を札のかなとした．
漢字で始まるボケは pykakasi で読みに変換したが，「辛い」（からい，つらい）のように複数の読みがありうるため，漢字で始まるボケは順守とみなした．
第1版の検証データのカルタ・穴埋め・お題文のボケ90件で，この規則の正解率は0.97であり，違反4件のうち2件を検出した．
VLM の教師の正解率は0.88〜0.98であり，規則と同程度であったが，速度は規則の数千分の一であった．
カルタ以外の形式で制約に反するボケは，ほとんどが宣伝や他のボケへの返信などの無関係な書き込みであった．

\subsection{穴埋めとお題文}
穴埋めとお題文の制約は，文字の一致では確かめられず，ボケが空欄や問いに意味の上で応じているかを判断する必要がある．
しかし，実際の違反はまれであり，検証データの穴埋め・お題文のボケ120件では1件しかなかった．
そこで，検証データの穴埋め・お題文のボケを，同じ形式の別のお題に付け替えた組120件を違反の代わりとし，元の組120件との判別で方法を比較した．

比較した方法は，VLM の教師（Qwen3-VL-32B，30B-A3B）にお題画像とボケを与えて問う方法，ボケの文の埋め込みとお題画像の埋め込みを連結して正しい組と付け替えた組を判別する MLP（学習用に2万件の組を作成），同じお題の他のボケとの類似度の3種類である．
類似度は，ボケ $ b $ の文の埋め込みを $ {\mathbf e}_b $ とし，同じお題の他のボケ（最大50件）のうち $ {\mathbf e}_b $ と最も似た5件の集合を $ {\mathbb N}_5(b) $ として，(\ref{eq_knn})で求めた．

\begin{equation}
    s(b) = \frac{1}{5} \sum_{b' \in {\mathbb N}_5(b)} \cos ( {\mathbf e}_b, {\mathbf e}_{b'} )
    \label{eq_knn}
\end{equation}

\noindent
文の埋め込みには ruri-v3-310m\cite{ref_ruri}の平均プーリングを用いた．
同じお題の他のボケは「お題が求める答えの方向」を表すため，お題の文を書き起こさなくても応答の妥当性を測れると考えた．
比較のため，最も似た5件ではなく他のボケ全ての埋め込みの平均とのコサイン類似度も求めた．

表\ref{tab_compliance}に結果を示す．
検出率は，正しい組の95\%を順守と判定するしきい値で，付け替えた組を検出できた割合である．

\begin{table}[t]
    \centering
    \caption{穴埋め・お題文の制約の順守の判別（正120組，付け替え120組）．}
    \label{tab_compliance}
    \small
    \begin{tabular}{lcccc}
        \toprule
        方法 & AUC & 検出率 & 穴埋め & お題文 \\
        \midrule
        教師 32B & 0.55 & 13\% & --- & --- \\
        教師 30B-A3B & 0.52 & 6\% & --- & --- \\
        組の分類器 & 0.57 & 10\% & 0.54 & 0.65 \\
        他のボケの平均 & 0.75 & 32\% & 0.66 & 0.83 \\
        類似度 $ s(b) $ & \textbf{0.81} & \textbf{39\%} & \textbf{0.73} & \textbf{0.88} \\
        \bottomrule
    \end{tabular}
\end{table}

VLM の教師は，付け替えた組もほぼ全て順守と判定し，画像の空欄や問いとボケの対応を確かめられなかった．
組の分類器も弱く，画像全体の埋め込みはお題の文の中身を表さないと考えられる．
同じお題の他のボケとの類似度が最も良く，特に最も似た5件との類似度は，典型から外れた正しいボケも拾える点で平均との類似度を上回った．
お題文では AUC 0.88 であったのに対し，穴埋めでは空欄に入る語がボケどうしで多様なため0.73にとどまった．

この得点は他のボケと似ているかを測るため，意外性のある正しいボケも低くなりうる．
そこで，しきい値で二分せず得点のまま付与し，生成の学習データから下位数\%を除く用途や，判定の特徴として用いることを提案する．
また，ボケの少ないお題では計算できないため，その場合にはお題の文の書き起こしが必要になる．
